\section{Methodology}
\label{sec:method}

\subsection{Edge-Weight Framework}

Let $G = (V, E)$ be the census-tract adjacency graph for a state,
where each vertex $v \in V$ represents one census tract and each edge
$(u,v) \in E$ represents a shared boundary.
The standard METIS minimum-edge-cut objective minimizes
$\sum_{(u,v) \in \partial P} w(u,v)$
where $\partial P$ is the set of edges crossing partition boundaries and
$w(u,v) = |\text{boundary}(u,v)|$ is the TIGER boundary length in metres.

We extend this with a composable edge-weight pipeline.
Each $\mathtt{EdgeWeighter}$ is a function $f : \mathbb{R}_{>0}^{|E|} \to \mathbb{R}_{>0}^{|E|}$
that transforms the current edge-weight map.
The final weights are the composition:
\[
  w^* = f_k \circ \cdots \circ f_2 \circ f_1(w^{\text{geo}})
\]
where $w^{\text{geo}}$ is the geographic base (TIGER boundary lengths).

\subsection{Subdivision Stickiness}

\begin{definition}[County stickiness weight]
For stickiness parameter $\alpha \geq 0$, the \emph{county-sticky} edge weighter is:
\[
  f_\alpha(u,v) = w(u,v) \cdot \bigl(1 + \alpha \cdot \mathbf{1}[\text{county}(u) = \text{county}(v)]\bigr).
\]
\end{definition}

When $\alpha = 0$, $f_0$ is the identity.
When $\alpha > 0$, intra-county edges cost $1 + \alpha$ times their geographic length,
while cross-county edges keep their original cost.
METIS therefore prefers to cut across county lines, which already represent
established administrative boundaries.

\paragraph{County derivation from GEOIDs.}
Census tract GEOIDs are 11-digit strings of the form $\text{SS CCC TTTTTT}$
where SS is the two-digit state FIPS code and CCC is the three-digit county FIPS
within the state.
The full county FIPS is the first five characters of the GEOID.
This information is stored in the TIGER \texttt{.adj.bin} adjacency file alongside
the adjacency graph, requiring no additional data download.

\begin{proposition}[Monotonicity, exact solver]
\label{prop:mono}
Let $P^*_\alpha$ denote the globally optimal $k$-way partition minimising
$\sum_{(u,v)\in\partial P} f_\alpha(u,v)$.
The number of county splits $\text{splits}(P^*_\alpha)$ is weakly decreasing in $\alpha$.
\end{proposition}

\begin{proof}
Let $P^*_\alpha$ be optimal under $f_\alpha$ and suppose it splits $s$ counties.
For any $\alpha' < \alpha$, we have $f_{\alpha'}(u,v) \leq f_\alpha(u,v)$ for all
$(u,v)$, so the cut cost of $P^*_\alpha$ under $f_{\alpha'}$ is at most its cost
under $f_\alpha$.
But the optimal partition under $f_{\alpha'}$ minimises cost further;
in particular, it can always route along county boundaries (cost unchanged
under $f_{\alpha'}$), so its county-split count is at most $s$.
By exchange argument, reducing $\alpha$ cannot force more county splits in an
optimal solution. \qed
\end{proof}

\begin{remark}
Proposition~\ref{prop:mono} applies to the \emph{exact} weighted minimum
$k$-way cut problem.
METIS~\citep{Karypis1998} is a multilevel heuristic and does not guarantee
globally optimal solutions; its output need not be monotone in $\alpha$ for
any fixed random seed.
Empirically, Section~\ref{sec:eval} shows that mean county splits over an
ensemble of seeds are weakly decreasing in $\alpha$ across all 44 multi-district
states, consistent with the monotone optimal-solution property.
\end{remark}

\subsection{Hierarchical Extension (Future Work)}
\label{sec:hierarchy}

The county weighter generalises to any level of the governmental subdivision
hierarchy.
Let $L = \{\text{county}, \text{MCD}, \text{place}, \text{VTD}\}$.
The hierarchical weighter formula is:
\[
  f_{\boldsymbol\alpha}(u,v) = w(u,v) \cdot \Bigl(1 + \sum_{\ell \in L} \alpha_\ell \cdot \mathbf{1}[\text{level}_\ell(u) = \text{level}_\ell(v)]\Bigr).
\]
We state this formula for completeness; \emph{the present paper validates and
recommends only $\alpha_{\text{county}}$}.
The other levels require additional TIGER downloads (COUSUB, PLACE, VTD files)
and spatial joins, and are left for future empirical validation.

\begin{remark}[Composition amplification]
When multiple levels are active simultaneously, a tract pair sharing county,
MCD, and place membership receives multiplier $(1 + \alpha_c + \alpha_m + \alpha_p)$.
For $\alpha = 5$ at each of three levels this is $16\times$ the base weight,
which may dominate the geographic signal.
Practitioners using multi-level composition should reduce per-level $\alpha$
proportionally (e.g., $\alpha = 2$ per level for three levels gives $7\times$,
comparable to county-only $\alpha = 5$).
We defer a formal analysis of multi-level balance behaviour to future work.
\end{remark}

\subsection{Composition with Other Signals}

The stickiness weighter composes with any other $\mathtt{EdgeWeighter}$ in the pipeline.
For example, combining county stickiness with the GeoSection split strategy
(B.8):
\[
  w^* = f_{\alpha_{\text{county}}} \circ f_{\text{geo}}(w^0)
\]
runs GeoSection using edge weights that penalise intra-county cuts.
The split strategy and the weight signal are orthogonal: adding $\alpha > 0$
does not change the GeoSection ratio-selection algorithm, only the weights it
optimises over.

\subsection{Implementation}

The $\mathtt{SubdivisionWeighter}$ is implemented in Rust as a \texttt{trait EdgeWeighter}
with a single \texttt{apply(EdgeMap) -> EdgeMap} method.
Concretely:
\begin{verbatim}
w'(u,v) = w(u,v) * (1 + alpha_county * same_county(u, v))
\end{verbatim}
where \texttt{same\_county(u,v)} is a precomputed boolean array derived from
GEOID prefixes.
The weighter is constructed once per state run and applied before partitioning.
Runtime overhead is $O(|E|)$, negligible relative to METIS partitioning.
The \texttt{--alpha-county} flag on \texttt{redist state} exposes this parameter
at the CLI.
